\clearpage
\fancyhead[L]{\small BURCKHARDT. TABLES DE LA LUNE.}
\SourceLeaf{558}
\begin{center}\small Göttingische gelehrte Anzeigen. Stück 39. Seite 388..392. 1813 März 8.\\[3pt]\rule{.72\linewidth}{.3pt}\end{center}

\emph{Tables astronomiques publiées par le bureau des longitudes de France. Tables de la lune, par} M. \textsc{Burckhardt}, \emph{membre de l'institut impérial, du bureau des longitudes de France et de plusieurs autres sociétés savantes. Décembre} 1812. Paris. Bei Madame Courcier. 88 Seiten in Quart.

Die Berechnung der Mondsörter hat durch die \textsc{Bürg}'schen Tafeln einen so hohen Grad von Genauigkeit erhalten, dass es ein gewagtes Unternehmen scheint, diese berühmten Tafeln noch übertreffen zu wollen. Inzwischen ist die möglichste Vollkommenheit der Mondstafeln in vielfacher Beziehung von so hoher Wichtigkeit, dass man allerdings einem so geschickten Astronomen, wie der Verfasser der vorliegenden Tafeln ist, für seine Bemühungen, diese Vollkommenheit noch zu erhöhen, den grössten Dank schuldig ist. Arbeiten dieser Art sind um so verdienstlicher, da ihnen nicht einige Monate, sondern \emph{Jahre} geopfert werden müssen, und sie, der Natur der Sache nach, nicht mehr durch \emph{glänzende} Erfolge belohnt werden können.

Der Verfasser hatte bei seiner Unternehmung einen doppelten Zweck. Zunächst nemlich wollte er den Tafeln eine etwas veränderte, bequemere und einfachere \emph{Form} geben: allein um diesen Zweck zu erreichen, begnügte er sich nicht damit, blos die Elemente der \textsc{Bürg}'schen Mondstafeln umzuschmelzen, sondern er gründete vielmehr die seinigen auf die eigene neue Bearbeitung von mehr als vier tausend Beobachtungen, so dass diese Tafeln als wahres alleiniges Eigenthum des Herrn \textsc{Burckhardt} anzusehen sind. Mit Recht konnte er hoffen, dass auf diese Weise die neuen Tafeln auch in Rücksicht auf \emph{Genauigkeit} noch einigen Vorzug erhalten würden, und in der That bestätigt dies die Vergleichung von drei hundert Beobachtungen, welche das Französische Büreau der Meereslänge sowohl mit den \textsc{Burckhardt}'schen, als mit den \textsc{Bürg}'schen Tafeln anstellen liess. 167 zu Greenwich und auf der kaiserlichen Sternwarte zu Paris beobachtete Längen gaben die Summe der Quadrate der Fehler nach den \textsc{Burckhardt}'schen Tafeln $= 4602''$, nach den \textsc{Bürg}'schen hingegen $= 7083''$; 137 andere in Paris auf der kaiserl. Sternwarte und auf der Militairschule beobachtete Längen gaben die Summe der Quadrate der Fehler nach \textsc{Burckhardt}'s Tafeln $4182''$, nach \textsc{Bürg}'s Tafeln $6439''$. Auch die Breiten stimmten, wie der Verfasser versichert, besser mit seinen eigenen Tafeln, als mit den \textsc{Bürg}'schen: die Grösse der Abweichung ist aber hier nicht angegeben. Wir hätten gewünscht, die Resultate dieser sämmtlichen Vergleichungen hier einzeln abgedruckt zu finden; theils wäre dadurch die Ueberzeugung von der hohen gegenwärtigen Vollkommenheit der Mondstafeln noch anschaulicher geworden, theils würde dadurch die künftige Prüfung, ob diese Vollkommenheit durch Hinzufügung einer oder der andern neuen Gleichung noch Etwas gewinnt, ungemein erleichtert sein. Ueberhaupt hätten wir in der Einleitung vor diesen Tafeln etwas mehr Ausführlichkeit gewünscht; wir sind zwar übrigens keineswegs für die weitläufigen Anweisungen, in welchen manche Verfasser von Tafeln allbekannte Dinge zum Ueberdruss wiederholen: aber das, was neuen Tafeln eigenthümlich ist, in der Kürze, aber doch vollständig und ausdrücklich, angezeigt zu finden, scheint uns doch ein billiger Wunsch, wenn es gleich nicht schwer ist, dies durch Analysirung der Tafeln selbst heraus zu finden.

Der vornehmste Unterschied der Form dieser Tafeln von derjenigen, welche seit \textsc{Tobias Meyer} von \textsc{Mason}, \textsc{Triesnecker} und \textsc{Bürg} beibehalten war, besteht darin, dass nicht die wahren, sondern die mittleren Sonnenlängen zum Grunde liegen. Auch Länge des Knoten und Perigeum werden hier nicht erst durch eine von der Sonnen-Anomalie abhängige Gleichung verbessert. Dagegen ist die Evection von den übrigen kleinen Gleichungen getrennt. Die Anzahl der kleinen Gleichungen, deren Argumente \SourceLeaf{559}alle, unabhängig von einander, unmittelbar aus der Tafel leicht entnommen werden, beträgt, die Nutation und zwei Störungsgleichungen von der Venus und dem Jupiter eingeschlossen, jetzt 32; hiermit wird das mittlere Argument der Evection verbessert; die Summe jener 32 Gleichungen und der Evection verbessert die Anomalie; die Summe der 32 Gleichungen, der Evection und der Mittelpunktsgleichung verbessert das Argument der Variation; und diese nebst den vorigen 34 Gleichungen, mit der mittlern Länge vereinigt, gibt die wahre Länge in der Bahn, die dann noch auf die Ekliptik reducirt wird. Bei den Argumenten der Evection, der Mittelpunktsgleichung, der Variation und bei der mittlern Länge muss zugleich noch die Säculargleichung zugezogen werden; mit der letztern ist die bekanntlich empirisch bestimmte kleine Ungleichheit von langer Periode in den Epochen für das 19. Jahrhundert vereinigt; bei andern Jahrhunderten muss man beide nach einer am Ende hinzugefügten besondern Tafel getrennt berechnen. In Rücksicht dieser kleinen Gleichung hält \textsc{Laplace} jetzt für das wahrscheinlichste, dass sie dem Cosinus der doppelten Länge des Mondsknoten, plus der Länge des Perigäum proportional sei, die Länge der Periode ist sonach 175 Jahre, der Coëfficient wird $= 12''5$ gesetzt, und ihr Ursprung liegt in einer vorausgesetzten Ungleichheit der nördlichen und südlichen Erd-Hemisphäre. Es ist zu wünschen, dass die Theorie in Beziehung auf diesen wichtigen Punkt noch mehr vervollkommnet werden möge. In den \textsc{Bürg}'schen Mondstafeln war diese Gleichung dem Sinus jenes Arguments weniger der dreifachen Länge des Sonnenperigäum proportional, und ihr Coëfficient $= 14''$ angenommen. Bei der Horizontal-Parallaxe hat sich Hr. \textsc{Burckhardt} ganz an die \textsc{Laplace}'sche Theorie gehalten; das Verhältniss derselben zum Horizontalhalbmesser des Mondes gründet sich auf die in den Vollmonden beobachtete Dauer der Durchgänge durch den Meridian. --- Am Schlusse des Werks sind noch ein Paar von Hrn. \textsc{Burckhardt} neu entwickelte Formeln angehängt, die mit Vortheil zur Berechnung der Neu- und Vollmonde angewandt werden können.
